\chapter*{Záver}  % chapter* je necislovana kapitola
\addcontentsline{toc}{chapter}{Záver} % rucne pridanie do obsahu
\markboth{Záver}{Záver} % vyriesenie hlaviciek

V rámci tejto práce sme vyvinuli webovú aplikáciu určenú pre kontrolu ekvivalentných úprav. Vyvíjali sme ju iteratívno-inkrementálnym prístupom. Najprv sme aplikáciu implementovali s minimálnou funkcionalitou a postupne ju vyvíjali pridávaním ďalších funkcionalít. Navrhli sme stav aplikácie a implementovali ho pomocou knižnice Redux. Postupne sme pridávali do aplikácie kontrolu rôznych ekvivalentných úprav a skolemizácie. Potrebovali sme pri tom rozšíriť dátový model pre reprezentáciu formúl vytvorený počas bakalárskej práce Jozefa Filipa. Taktiež sme rozšírili nástroj na syntaktickú analýzu vytvorený počas bakalárskej práce Milana Cifry.

Nakoniec sme aplikáciu otestovali počas výučby predmetu Logika pre informatikov. Študenti po práci s aplikáciu vyplnili dotazník, podľa ktorého sme zhodnotili, že študenti ocenili jednoduchosť aplikácie. Tiež sme zhodnotili, že používateľnosť aplikácie je dobrá. Výsledná verzia aplikácie je vhodná na použitie pri výučbe predmetu.

Existuje stále priestor na zlepšenie v viacerých oblastiach. Najčastejšou pripomienkou študentov boli chybové hlášky pri nesprávnej úprave. Jedným zlepšením by preto mohlo byť ich vylepšenie, poprípade skonkretizovanie. Ďalším možným zlepšením by mohla byť samotná kontrola úprav. V tejto chvíli aplikácia vie kontrolovať iba jeden typ úpravy. Vylepšením by mohlo byť to, že niektoré základné úpravy (ako napríklad asociativita a komutativita) by sa kontrolovali aj pri iných úpravách.

Vďaka tejto práci, sme mali možnosť zoznámiť sa s tvorbou moderných webových aplikácií. Práca nás naučila o knižniciach React a Redux, ktoré boli pre nás nové. Zoznámili sme sa so správou stavu webových aplikácii. Tiež sme nadobudli cenné skúsenosti s vývojom aplikácie a práce správou verzií kódu pomocou platformy GitHub.  